%GIÁO SƯ PHAN ĐÌNH DIỆU (1936-2018): KHÍ PHÁCH \& TRÍ TUỆ (Trương Huy San)
 
Nguồn: \href{http://nghiencuuquocte.org/forums/topic/giao-su-phan-dinh-dieu-1936-2018-khi-phach-va-tri-tue/}{Báo Nghiên cứu quốc tế} 
và \href{https://www.facebook.com/Osinhuyduc/posts/1620411811327327}{Facebook của tác giả} 




Sáng nay tôi hỏi 5 bạn sinh trong thập niên 1980s, 3 bạn nói không biết GS Phan Đình Diệu. Tôi buồn nhưng không bất ngờ. Ông là một người mà thế hệ chúng tôi ngưỡng mộ, cả về tài năng và sự chính trực. Có lẽ vì sự chính trực ấy mà ông rất ít khi có mặt trên những bục vinh quang. Nhưng, những đóng góp thầm lặng của ông đặc biệt là trong vai trò đưa công nghệ thông tin vào VN đang ảnh hưởng lên cả những thế hệ không còn biết ông là ai nữa.
 

 
GS Phan Đình Diệu là một trong những người đầu tiên gầy dựng ngành khoa học tính toán của miền Bắc (1968). Chiếc máy vi tính đầu tiên, có thể coi là của Châu Á, ra đời năm 1980 [được thiết kế với chip Intel 8080A nên được đặt tên là VT80] là từ phòng thí nghiệm Viện Tin học của ông, bắt đầu từ những hợp tác trước đó với Viện Kỹ thuật Quân sự có sự giúp sức cả về chuyên môn và vật chất của ông André Trương Trọng Thi, người được coi là cha đẻ của máy tính cá nhân.

Theo tiến sỹ Vũ Duy Mẫn: “Năm 1981, sau khi lắp ráp thành công máy vi tính VT81 và cài đặt được ngôn ngữ lập trình Basic, Viện quyết định thử nghiệm đưa máy vi tính ứng dụng vào quản lý xí nghiệp cỡ nhỏ và vừa ở Xí nghiệp máy may Sinco”.

Nhưng, cho dù có nhiều nỗ lực, ngành công nghiệp sản xuất máy tính chỉ dừng lại ở đó. Tin học là một ngành không thể phát triển trong một quốc gia tự đóng cửa về chính trị và bị cô lập về thương mại với thế giới bên ngoài. Rất tiếc là khi tình hình trong nước bắt đầu cởi mở hơn thì GS Nguyễn Văn Hiệu được đưa lên thay giáo sư Trần Đại Nghĩa, làm Viện trưởng Viện khoa học Việt Nam. Năm 1993, ông Hiệu loại GS Phan Đình Diệu khi ngành tin học Việt Nam cần một người lãnh đạo có tâm và có tầm nhất.
GS Hiệu năm ấy 45 tuổi. Ông có mọi ưu đãi của Chế độ, từ kinh phí nghiên cứu tới quyền hạn. Nắm Viện, ông Hiệu kêu gọi “các anh các chị hãy tự cứu mình”. Các công ty được thành lập, Viện khoa học Việt Nam nhanh chóng trở thành một nơi “nửa hàn lâm, nửa chợ trời”[theo TS Giang Minh Thế]. Hàng nghìn cán bộ tài năng trôi giạt khắp nơi. Nền tin học Việt Nam chuyển một cách dứt khoát sang nghề… buôn máy tính.
Dù GS Phan Đình Diệu là người viết đề cương thành lập viện Công nghệ thông tin theo mô hình mới và khi bỏ phiếu thăm dò, ông nhận được tín nhiệm của đa số tuyệt đối, GS Nguyễn Văn Hiệu vẫn chọn một người có số phiếu thấp hơn vì lý do “biết làm kinh tế”. GS Phan Đình Diệu tuyên bố từ chức Viện phó Viện Khoa học Việt Nam và ra khỏi biên chế, những chuyên gia trẻ hết lòng vì khoa học như Nguyễn Chí Công, Vũ Duy Mẫn, Lê Hải Khôi, Giang Công Thế… bắt đầu phiêu bạt.

Cho dù vậy, theo GS Chu Hảo – thứ trưởng Bộ Khoa học Công nghệ thời kỳ internet được đưa vào VN – “GS Phan Đình Diệu vẫn là linh hồn của các chính sách phát triển công nghệ thông tin của nước ta”. Ông là Chủ tịch hội Tin học và là Phó Chủ tịch Uỷ ban Quốc gia phát triển Công nghệ Thông tin cho tới năm 1998.
Nếu như, GS Nguyễn Văn Hiệu nổi tiếng là một người biết sử dụng khoa học để làm chính trị và leo lên tới các đỉnh cao danh vọng thì GS Phan Đình Diệu lại là một con người mà ngay cả trong chính trị cũng chỉ tham gia với tinh thần khoa học.

Theo ông Trần Việt Phương – thư ký riêng của Thủ tướng Phạm Văn Đồng – từ đầu thập niên 1960s, khi được Thủ tướng Phạm Văn Đồng hỏi, GS Phan Đình Diệu đã cho rằng, hệ thống kinh tế Xô Viết sẽ sụp đổ nếu tiếp tục vận hành như thế này[Hy vọng là Trung tâm Lưu trữ quốc gia vẫn còn giữ được ý kiến này của GS Phan Đình Diệu].

Rất có thể, quyết định không bổ nhiệm GS Phan Đình Diệu vào năm 1993 của ông Hiệu còn có “lý do chính trị”. Từ năm 1989, GS Phan Đình Diệu có nhiều phát biểu về “đa nguyên” và đặc biệt là các ý kiến thẳng thắn của ông góp ý cho Hiến pháp 1992.

Tháng 10-1988, sau khi giải quyết xong vấn đề “đa đảng”[giải tán hai đảng Dân chủ và Xã hội], Tổng bí thư Nguyễn Văn Linh bắt đầu “chống đa nguyên” nhắm vào ông Trần Xuân Bách. Tháng 3-1989, Hội nghị Trung ương 6 khẳng định “Không chấp nhận chủ nghĩa đa nguyên”. Ngay sau đó trên báo Đảng, cả Trưởng ban Tuyên huấn Trung ương Trần Trọng Tân lẫn “nhà nghiên cứu” Trần Bạch Đằng đều có bài, trên tinh thần “Mác Xít”, nhấn mạnh rằng, “Nếu hiểu đa nguyên theo nghĩa triết học thì từ lâu đã bị phê phán là sai lầm. Chúng ta không chấp nhận”.

Ngày 15-8-1989, GS Phan Đình Diệu có bài trên Sài Gòn Giải Phóng, đáp trả các luận điệu trên đây. Ông cho rằng, muốn “tìm đường đi cho đất nước trong thế giới hiện đại” thì phải “vận dụng trí tuệ của thời đại” từ “nhiều nguồn tri thức”. Theo ông, “Những gì xảy ra trong thế kỷ hai mươi là điều không thể hình dung đối với những bộ óc, dù là vĩ đại, của giữa thế kỷ mười chín”.

Sau khi đọc bài của Giáo sư Phan Đình Diệu, từ Hà Nội, Bí thư Trần Xuân Bách gửi cho TBT Sài Gòn Giải Phóng Tô Hoà một tấm danh thiếp, mặt sau ghi: “Chuyển giùm anh Phan Đình Diệu, tôi ca ngợi bài này”. Đây là bài báo cuối cùng mà Tô Hoà cho đăng với tư cách tổng biên tập Sài Gòn Giải Phóng. Ngày 28-3-1990, Trung ương cách chức ông Trần Xuân Bách.

Hai năm sau, ngày 12-3-1992, trong một phiên họp góp ý cho Hiến pháp 1992, GS Phan Đình Diệu đã có một phát biểu rất gây tiếng vang ở Uỷ ban Trung ương MTTQ Việt Nam. Ông đề nghị “tạm gác lại” việc đưa vào Hiến pháp các thuật ngữ “chủ nghĩa xã hội, chủ nghĩa Marx – Lenin”. Ông đề nghị phải dũng cảm để thấy rằng, “mô hình CNXH, cũng như Marx – Lenin không còn đáp ứng được mục tiêu phát triển đất nước trong giai đoạn hiện nay”. Giai đoạn xây dựng một “tổ quốc Việt Nam trên tinh thần đoàn kết, hoà hợp và tự cường”.
Theo GS Phan Đình Diệu thì Việt Nam dứt khoát phải có “nhà nước pháp quyền”, chứ không thể tiếp tục “chuyên chính vô sản” hay “pháp quyền nửa vời”. Đặc biệt, ý kiến của ông về Mặt trận và các đoàn thể cho tới ngày nay vẫn còn rất thiết thực. Ông đề nghị, Hiến pháp chỉ cần coi lập hội là quyền tự do của người dân; từ bỏ việc nhà nước hoá các đoàn thể. Ông kêu gọi các đoàn thể thay vì dựa dẫm vào nhà nước phải lấy quần chúng làm gốc rễ, các cơ quan đảng, đoàn thể không ăn lương từ ngân sách.

GS Phan Đình Diệu không sử dụng tài năng và danh tiếng của mình để mưu cầu bổng lộc hay địa vị. Những năm sau 1975, trong những ngày đói kém, từ Paris trở về, hành lý của ông vẫn chỉ có sách và các “bản mạch” dùng cho sản xuất máy tính. Khi phải xuất hiện trên các diễn đàn thì, cho dù ở giai đoạn nào, ý kiến của ông vẫn là ý kiến của một con người yêu nước, khí phách và trí tuệ.